\documentclass[11pt,reqno]{amsart}
\usepackage{geometry}                % See geometry.pdf to learn the layout options. There are lots.
\geometry{letterpaper}                   % ... or a4paper or a5paper or ... 
%\geometry{landscape}                % Activate for for rotated page geometry
%\usepackage[parfill]{parskip}    % Activate to begin paragraphs with an empty line rather than an indent
\usepackage{graphicx}
\usepackage{amssymb}
\usepackage{epstopdf}
\DeclareGraphicsRule{.tif}{png}{.png}{`convert #1 `dirname #1`/`basename #1 .tif`.png}

\usepackage{color}
\definecolor{darkblue}{rgb}{0.0,0.0,0.6}
\usepackage[pdfauthor={Tsogtgerel Gantumur},pdftitle={Assignment},pdfsubject={PDE},colorlinks,citecolor=darkblue,linkcolor=black,urlcolor=darkblue]{hyperref}

\usepackage{enumerate}
\usepackage{nicefrac}

\usepackage{mathrsfs}

\newcommand{\R}{\mathbb{R}}
\newcommand{\N}{\mathbb{N}}
\newcommand{\tstD}{\mathscr{D}}
\newcommand{\tstC}{\mathscr{C}}
\newcommand{\semP}{\mathscr{P}}
\newcommand{\eps}{\varepsilon}
\newcommand{\supp}{\mathrm{supp}}
\newcommand{\exd}{\mathrm{d}}

\title{Math 580 Assignment 1}
\author{Due Monday September 23}
\date{Fall 2013}                                           % Activate to display a given date or no date

\begin{document}
\maketitle

\begin{enumerate}[1.]
\item
Exercises 2, 4, 6, 7, 11 from the lecture notes \href{/legacy/courses/math580f13/harmonic.pdf}{\em Harmonic functions}.
\item
Prove that the space of harmonic functions on an open set $\Omega\subset\R^n$ ($n\geq2$) is infinite dimensional.
\item
Let $\Omega\subset\R^n$ be a bounded domain,
and consider the boundary value problem
\begin{equation}\label{e:nondir}\tag{$*$}
\Delta u = f(u)\quad\textrm{in }\Omega,
\qquad
u = 1\quad\textrm{on }\partial\Omega.
\end{equation}
Prove the followings.
\begin{enumerate}
\item[a)]
Any solution of \eqref{e:nondir} in $C^2(\Omega)\cap C(\bar\Omega)$ with $f(u)=u^m$ where $m\in\N$ is odd, must satisfy $0\leq u\leq 1$ in $\bar\Omega$, and is unique.
\item[b)]
The only solution of \eqref{e:nondir} in $C^2(\Omega)\cap C(\bar\Omega)$ with $f(u)=u-u^{-1}$ is $u\equiv1$.
\end{enumerate}
\item
Let $\phi:(0,1]\to\R$ be a continuous nonnegative function satisfying
$$
\int_\eps^1\phi(r)r^{n-1}\exd r \leq M,
$$
for any $\eps>0$, with a fixed constant $M$.
Define $f(x)=\phi(|x|)$ for $x\in B_1$ where $B_1=\{x\in\R^n:|x|<1\}$.
Prove in complete detail that $f\in L^1(B_1)$.
\item
Let $u\in\tstC^2(\R^n)$ be an entire harmonic function in $\R^n$.
Prove the followings.
\begin{enumerate}[(a)]
\item
If $u\in L^p(\R^n)$ for some $1\leq p<\infty$ then $u\equiv0$.
\item
If $u$ satisfies $u(x)\geq-C(1+|x|)^m$ for some constants $C$ and $m\in\mathbb{N}$,
then $u$ is a polynomial of degree less or equal to $m$.
\item
Any tangent hyperplane to the graph of $u$ intersects the graph more than once.
\end{enumerate}
\end{enumerate}

\section*{Homework policy}

You are welcome to consult each other provided (1) you list all people and sources who aided you, or whom you aided and (2) you write-up the solutions independently, in your own language. If you seek help from other people, you should be seeking general advice, not specific solutions, and must disclose this help.
This applies especially to internet fora such as \texttt{MathStackExchange}.

Similarly, if you consult books and papers outside your notes, you should be looking for better understanding of or different points of view on the material, not solutions to the problems.
\end{document}  











